% Generated draft: requires mathematical and source review.
\documentclass[11pt]{article}
\usepackage[margin=25mm]{geometry}
\usepackage{fontspec}
\setmainfont{Latin Modern Roman}
\usepackage{amsmath,amssymb,mathtools,unicode-math}
\setmathfont{Latin Modern Math}
\usepackage{enumitem,needspace,xurl,hyperref}
\setlength{\parindent}{0pt}
\setlength{\parskip}{5pt}
\setlength{\emergencystretch}{3em}
\newcommand{\N}{\mathbb{N}}
\newcommand{\Z}{\mathbb{Z}}
\newcommand{\Q}{\mathbb{Q}}
\newcommand{\R}{\mathbb{R}}
\newcommand{\C}{\mathbb{C}}
\newcommand{\F}{\mathbb{F}}
\newcommand{\A}{\mathbb{A}}
\newcommand{\PP}{\mathbb{P}}
\newcommand{\E}{\mathbb{E}}
\newcommand{\eps}{\varepsilon}
\newcommand{\dd}{\,\mathrm{d}}
\begin{document}
\section*{UnsolvedMath ID: 30001231}
\textbf{Source problem code:} catalogue entry
\subsection{Definitions and mathematical statement}
A standard one-point Seshadri constant is defined in the following conditional setting. Let \(X\) be an integral projective variety over an algebraically closed field \(k\), let \(L\) be a nef line bundle on \(X\), and let \(x\in X\) be a closed point. Define
\[
\varepsilon(L;x)=\inf_{\substack{C\subseteq X\text{ an integral curve}\\x\in C}}\frac{\deg(L|_C)}{\operatorname{mult}_x C},
\]
where \(\operatorname{mult}_x C\) denotes the multiplicity of \(C\) at \(x\). Here nef means that \(\deg(L|_C)\geq 0\) for every integral curve \(C\subseteq X\). These conventions supply a standard definition; the source does not specify that this is its exclusive intended setting.

The supplied research question asks what geometric information Seshadri constants encode and seeks upper bounds, lower bounds, or explicit values on projective surfaces and on higher-dimensional projective varieties. Under the one-point interpretation above, this asks for geometric consequences of \(\varepsilon(L;x)\) and estimates or computations of this number for specified data \((X,L,x)\), with \(\dim X=2\) or \(\dim X\geq 3\).

The record does not specify a ground field, smoothness hypotheses, whether \(L\) is ample or merely nef, the points under consideration, the permitted dependence of bounds, or whether one-point, multipoint, or global Seshadri constants are intended. Consequently, it supplies a research programme rather than a uniquely quantified conjecture. The conditional definition above must not be read as resolving these missing qualifications.
\subsection{Short English statement}
Seshadri constants measure local positivity of a line bundle through its degrees on curves relative to their multiplicities at a point. The supplied question asks which geometric features they detect and how to bound or compute them on surfaces and higher-dimensional projective varieties. The intended hypotheses and precise targets remain unspecified.
\subsection{Sources}
\begin{itemize}
\item Supplied local UnsolvedMath registry record, permanent ID \(30001231\), title: \emph{Geometric Information Encoded by Seshadri Constants}; exact source problem code: \texttt{catalogue entry}. The record gives proposed year 2009 and names no individual proposer.
\item \emph{Multiplier Ideal Sheaves in Algebraic and Complex Geometry} (2009), Oberwolfach Reports, pp.\ 1101--1156, DOI \texttt{10.4171/owr/2009/21}. Supplied source URL and literature-source link labelled ``original report'': \url{https://doi.org/10.4171/owr/2009/21}.
\item Literature assessment dated \texttt{2026-08-21}: the supplied curation labels the problem open and states that these extracted guiding questions about geometric information and effective values of Seshadri constants remain open in this generality. This is supplied dated curation, not an independently verified literature assessment. The background explicitly requests independent verification of time-sensitive status and cited claims.
\item The registry marks the statement as exact and reports manual verification of its coherence and freedom from extraction artifacts. This supplied verification does not resolve the missing mathematical scope identified above. No problem-specific page, named speaker, or separate authoritative citation for the definition is supplied.
\end{itemize}
\end{document}
